Published values are rarely commensurable, so no side-by-side display of results is attempted;
Table~\ref{tab:famreq} states what each family requires and how far this corpus has taken it. Decoding
protocols overlap most, and only within a modality and a shared public label schema. Otherwise probes
report AUROC on different concept sets, splits, label definitions, decoder capacities and controls, which
obscures differences of a few points. Discovery overlaps least: width, sparsity, layer and
interpretation all vary, cross-seed stability is seldom reported, and nominal concepts cannot be
matched. Intervention records can be inspected one by one, but their end-points, magnitudes and controls
differ, so effect sizes cannot be compared. Even the most compatible stratum, six chest-radiograph
studies using linear decoders on CheXpert or MIMIC-CXR, differs in targets, label versions, splits and
regularisation, and so shows six distinct estimands rather than a poolable set. Independent replication
by a second group appears to be rare, and the same-data medical--general comparisons confound domain with
data, objective, architecture and protocol (\S\ref{sec:findings}).

\textbf{Modality and family.} Pathology appears in \Npathstud{} of the \Ncore{} studies and chest
radiography in \Ncxrstud{} (\Npathrec{} and \Ncxrrec{} records); every other modality is markedly rarer
(Fig.~\ref{fig:corpus}c). Chest radiography supplies \Nintcxr{} of the \Nintrec{} intervention records
and \Ntracecxr{} of the \Ntracerec{} tracing records, so the interpretability of medical multimodal
LLMs is largely a chest-radiograph literature. Pathology dominates geometry and discovery and enters
intervention only through post-processing of frozen embeddings. Modality is retained in the evidence
tables because a claim rarely transfers across modalities.

\textbf{Study quality.} The quality fields are the subset of the imaging-AI and clinical-AI reporting
checklists~\cite{mongan2020claim,norgeot2020minimum} on which a representation-level claim can be scored
from the paper; they record what was reported, not whether it was done well. Of the \Ncore{} core
studies, \Npreprint{} are preprints, \Nconf{} conference papers, \Nwork{} workshop papers and \Njour{}
journal articles; \Ngrey{} of the preprints are grey-literature deposits and \Nsingle{} are
single-author. Excluding the grey deposits lowers the corpus-breadth tier of F8a$''$ alone and changes
no direction; excluding all \Ngreysingle{} lowers F3c as well. Of these studies, \Nextyes{} evaluate on a cohort or
dataset unused for probe, dictionary or intervention training, \Nextno{} do not and \Nextunk{} are
indeterminate; patient-, slide- or subject-level splits are stated in \Nsplityes{} and unstated in
\Nsplitno{}; seeds, repeated runs, intervals or tests accompany the main result in \Nseedyes{} and not
in \Nseedno{}, with \Nseedunk{} indeterminate. Of the \Nclaims{} records, \Ndirpos{} report a positive
result, \Ndirmix{} a mixed or partial one and \Ndirnull{} a null or negative one. Null results this rare
call for caution about publication and reporting bias in every statement of \S\ref{sec:findings}, and a
corpus that relies more on internal than on external validation is a second reason to describe breadth
rather than to pool effect sizes.

\textbf{Concept provenance.} Among the records, \Nprovlabel{} use structured schemas, \Nprovtext{}
text-mined concepts, \Nprovexpert{} expert vocabularies, \Nprovmodel{} model-generated banks and
\Nprovnone{} no prior concept, \Nposthocname{} of the last naming features post hoc. Expert-defined and
model-generated concepts occur almost only in intrinsic designs (\S\ref{sec:intrinsic}). The
best-supported decoding findings (F1a and F2 in \S\ref{sec:findings}) rest on labels automatically
extracted from chest-radiograph reports~\cite{irvin2019chexpert,johnson2019mimicjpg}, so their concept
validity is that of a labeller rather than of clinician judgement (\S\ref{sec:eval-data}). Discovery
without a prior concept shifts the burden to post-hoc naming, which fewer than half of the dictionary
studies validate with clinicians (\S\ref{sec:eval-what}). Provenance thus locates where a claim may fail
for reasons unrelated to the model.
